\documentclass[11pt]{letter}

\usepackage{geometry}
\geometry{margin=1in}
\usepackage{setspace}
\usepackage{hyperref}

\signature{Decory Edwards}
\address{Decory Edwards \\ Johns Hopkins University \\ Department of Economics}

\begin{document}

\begin{letter}{Hiring Committee \\ PDT Partners \\ New York, NY}

\opening{Dear Hiring Committee,}

I am writing to apply for the Quantitative Researcher position at PDT Partners. I am currently a PhD candidate in Economics at Johns Hopkins University. My training combines mathematical reasoning, computational modeling, and data driven empirical analysis. I am motivated by problems that require precise logic, careful implementation, and constant validation against data. PDT’s research driven approach and emphasis on rigor align closely with how I approach research and collaboration.

My research agenda is at the intersection of wealth inequality and macroeconomics. In my job market paper, I build and estimate a structural heterogeneous agent model with returns that differ across households. The core contribution is to show how heterogeneity in returns and income risk jointly shape the wealth distribution and consumption behavior. Relative to closely related work, my framework performs better at matching the lower and middle portions of the wealth distribution. This difference matters quantitatively. Models that focus primarily on returns heterogeneity can fit the top of the wealth distribution closely but tend to generate too much wealth among the bottom half of households. In my framework, income uncertainty generates a precautionary saving motive that produces genuinely low wealth households. This leads to a realistic distribution of marginal propensities to consume and aggregate consumption responses that align with empirical estimates.

A key feature of my research is the heavy use of computation to solve and simulate models that cannot be handled analytically. I write code to solve dynamic programming problems, simulate outcomes under stochastic environments, and evaluate model performance across specifications. I work with large and complex datasets and place a strong emphasis on robustness and interpretability. This work requires disciplined problem solving and careful validation, skills that map naturally to developing and improving proprietary trading models.

I am particularly interested in PDT Partners because of its small team structure and its focus on research quality over scale. The opportunity to work closely with senior researchers while contributing meaningfully to live research projects is especially appealing. I value environments where merit and contribution drive discussion and where ideas are evaluated through evidence and performance.

My economics training has provided a strong foundation in probability, optimization, and statistical inference. I am comfortable translating theory into working code and conducting independent research while collaborating closely with others. I have experience programming in Python and related tools and enjoy working at whiteboards to refine ideas and models. These skills are central to my research process and would be directly applicable in this role.

I am also enthusiastic about the prospect of living and working in New York City and participating in a hybrid research environment. I value in person collaboration while also appreciating focused research time. I see this role as an opportunity to apply my academic training in a setting where research has immediate real world impact.

I would welcome the opportunity to contribute my skills and perspective to PDT Partners. Thank you for your time and consideration.

\closing{Sincerely,}

\end{letter}
\end{document}
